%%%%%%%%%%%%%%%%%%%%%%%%%%%%%%%%%%%%%%%%%
% Presentación Beamer - Plantilla LaTeX
% Versión 2.0 (8 de marzo de 2022)
% Plantilla original: https://www.LaTeXTemplates.com
% Autor: Vel (vel@latextemplates.com)
% Licencia: CC BY-NC-SA 4.0

% Este modelo de presentación fue 
% creado a partir del modelo de Giovanni Spadaro.
% Disponible en: https://github.com/Giovo17/presentation-template-unict-lm-data
%
% Adaptado por Lautaro Lombardi para crear una versión especial 
% para los estudiantes de Facultad de Matemática, Astronomía, Física y Computación (FAMAF - UNC).
% Disponible en: https://github.com/lautalom/beamer_FAMAF
%%%%%%%%%%%%%%%%%%%%%%%%%%%%%%%%%%%%%%%%%

%----------------------------------------------------------------------------------------
% CLASE DEL DOCUMENTO Y CONFIGURACIONES BÁSICAS
%----------------------------------------------------------------------------------------
\documentclass[
    11pt,               % Tamaño estándar de la fuente
    % t,                % Alinear verticalmente en la parte superior
    %aspectratio=169,   % Definir proporción 16:9
]{beamer}
\graphicspath{{img/}}         % Define el directorio de las imágenes

%---------------------------------------------------------------------------------------
% PAQUETES NECESARIOS
%----------------------------------------------------------------------------------------
\usepackage{
    booktabs,     % Mejora la apariencia de las líneas en tablas
    palatino,     % Define Palatino como fuente principal
    subcaption    % Soporte para subfiguras
}
\usepackage[default]{opensans}  % Define Open Sans como fuente secundaria
%----------------------------------------------------------------------------------------
% Paquetes y configuraciones para código
%----------------------------------------------------------------------------------------
% Paquetes necesarios
\usepackage[utf8]{inputenc}
\usepackage{listings}
\usepackage{xcolor}

% Colores de Visual Studio Code Light Theme
\definecolor{vscBackground}{RGB}{255,255,255}    % Fondo blanco
\definecolor{vscKeyword}{RGB}{175,0,219}         % Púrpura para palabras clave
\definecolor{vscString}{RGB}{163,21,21}          % Rojo para cadenas
\definecolor{vscComment}{RGB}{0,128,0}           % Verde para comentarios
\definecolor{vscFunction}{RGB}{121,94,38}        % Marrón para funciones
\definecolor{vscNumber}{RGB}{9,134,88}           % Verde oscuro para números
\definecolor{vscOperator}{RGB}{175,0,219}        % Púrpura para operadores
\definecolor{vscText}{RGB}{0,0,0}                % Texto negro
\definecolor{vscLineNr}{RGB}{128,128,128}        % Gris para números de línea

% Configuração geral do listings para UTF-8
\lstset{
    inputencoding=utf8,
    extendedchars=true,
    literate=%
        {á}{{\'a}}1 {é}{{\'e}}1 {í}{{\'i}}1 {ó}{{\'o}}1 {ú}{{\'u}}1
        {Á}{{\'A}}1 {É}{{\'E}}1 {Í}{{\'I}}1 {Ó}{{\'O}}1 {Ú}{{\'U}}1
        {à}{{\`a}}1 {è}{{\`e}}1 {ì}{{\`i}}1 {ò}{{\`o}}1 {ù}{{\`u}}1
        {À}{{\`A}}1 {È}{{\'E}}1 {Ì}{{\`I}}1 {Ò}{{\`O}}1 {Ù}{{\`U}}1
        {ã}{{\~a}}1 {ẽ}{{\~e}}1 {ĩ}{{\~i}}1 {õ}{{\~o}}1 {ũ}{{\~u}}1
        {Ã}{{\~A}}1 {Ẽ}{{\~E}}1 {Ĩ}{{\~I}}1 {Õ}{{\~O}}1 {Ũ}{{\~U}}1
        {â}{{\^a}}1 {ê}{{\^e}}1 {î}{{\^i}}1 {ô}{{\^o}}1 {û}{{\^u}}1
        {Â}{{\^A}}1 {Ê}{{\^E}}1 {Î}{{\^I}}1 {Ô}{{\^O}}1 {Û}{{\^U}}1
        {ç}{{\c c}}1 {Ç}{{\c C}}1
        {º}{{\textordmasculine}}1
        {ª}{{\textordfeminine}}1
}

% Configurações base comum para todas as linguagens
\lstdefinestyle{baseStyle}{
    backgroundcolor=\color{vscBackground},
    basicstyle=\ttfamily\small\color{vscText},
    breakatwhitespace=false,
    breaklines=true,
    captionpos=b,
    keepspaces=true,
    numbers=left,
    numbersep=5pt,
    showspaces=false,
    showstringspaces=false,
    showtabs=false,
    tabsize=4,
    frame=single,
    framerule=0.8pt,
    rulecolor=\color{gray!20},
    numberstyle=\tiny\color{vscLineNr},
    keywordstyle=\color{vscKeyword},
    commentstyle=\color{vscComment}\itshape,
    stringstyle=\color{vscString},
    emphstyle=\color{vscFunction},
    columns=flexible,
    basewidth={0.5em,0.45em},
    inputencoding=utf8,
    extendedchars=true
}

%----------------------------------------------------------------------------------------
% Python
%----------------------------------------------------------------------------------------
\lstdefinestyle{pythonStyle}{
    style=baseStyle,
    language=Python,
    morekeywords={self,None,True,False,import,from,as,def,class,return,yield,
                  for,while,if,else,elif,try,except,finally,with,lambda,
                  async,await,break,continue,global,nonlocal,pass,raise},
    morekeywords=[2]{print,len,range,type,int,str,float,list,dict,set,
                     tuple,max,min,sum,sorted,enumerate,zip,map,filter,
                     any,all,abs,round,pow,divmod},
    keywordstyle=[2]\color{vscFunction},
    sensitive=true
}

\lstnewenvironment{python}[1][]{\lstset{style=pythonStyle, #1}}{}
\newcommand{\pyinline}[1]{\lstinline[style=pythonStyle]!#1!}
\newcommand{\inputpython}[2][]{\lstinputlisting[style=pythonStyle,#1]{#2}}

%----------------------------------------------------------------------------------------
% C Language
%----------------------------------------------------------------------------------------
\lstdefinestyle{cStyle}{
    style=baseStyle,
    language=C,
    morekeywords={include,define,void,int,char,float,double,long,unsigned,
                  struct,union,enum,typedef,const,static,extern,register,
                  auto,volatile,sizeof,return,if,else,for,while,do,switch,
                  case,break,continue,default,goto},
    morekeywords=[2]{printf,scanf,malloc,free,calloc,realloc,fopen,fclose,
                     fprintf,fscanf,strcpy,strlen,strcat},
    keywordstyle=[2]\color{vscFunction},
    sensitive=true
}

\lstnewenvironment{clang}[1][]{\lstset{style=cStyle, #1}}{}
\newcommand{\clinline}[1]{\lstinline[style=cStyle]!#1!}
\newcommand{\inputclang}[2][]{\lstinputlisting[style=cStyle,#1]{#2}}

%----------------------------------------------------------------------------------------
% C++
%----------------------------------------------------------------------------------------
\lstdefinestyle{cppStyle}{
    style=baseStyle,
    language=C++,
    morekeywords={class,private,protected,public,template,typename,namespace,
                  using,new,delete,this,friend,virtual,override,final,explicit,
                  mutable,constexpr,nullptr,noexcept,static_cast,dynamic_cast,
                  const_cast},
    morekeywords=[2]{cout,cin,endl,vector,string,map,set,queue,stack,pair,
                     begin,end,push_back,pop_back,emplace_back,size,empty},
    keywordstyle=[2]\color{vscFunction},
    sensitive=true
}

\lstnewenvironment{cpp}[1][]{\lstset{style=cppStyle, #1}}{}
\newcommand{\cppinline}[1]{\lstinline[style=cppStyle]!#1!}
\newcommand{\inputcpp}[2][]{\lstinputlisting[style=cppStyle,#1]{#2}}

%----------------------------------------------------------------------------------------
% R Language
%----------------------------------------------------------------------------------------
\lstdefinestyle{rStyle}{
    style=baseStyle,
    language=R,
    morekeywords={if,else,repeat,while,function,for,in,next,break,TRUE,FALSE,
                  NULL,Inf,NaN,NA,NA_integer_,NA_real_,NA_complex_,NA_character_},
    morekeywords=[2]{library,require,attach,detach,source,setwd,options,
                     data.frame,read.csv,write.csv,list,matrix,array},
    keywordstyle=[2]\color{vscFunction},
    sensitive=true
}

\lstnewenvironment{rlang}[1][]{\lstset{style=rStyle, #1}}{}
\newcommand{\rlinline}[1]{\lstinline[style=rStyle]!#1!}
\newcommand{\inputrlang}[2][]{\lstinputlisting[style=rStyle,#1]{#2}}

%----------------------------------------------------------------------------------------
% Java
%----------------------------------------------------------------------------------------
\lstdefinestyle{javaStyle}{
    style=baseStyle,
    language=Java,
    morekeywords={abstract,assert,boolean,break,byte,case,catch,char,class,
                  const,continue,default,do,double,else,enum,extends,final,
                  finally,float,for,if,implements,import,instanceof,int,
                  interface,long,native,new,package,private,protected,public,
                  return,short,static,strictfp,super,switch,synchronized,this,
                  throw,throws,transient,try,void,volatile,while},
    morekeywords=[2]{String,System,out,println,printStackTrace,ArrayList,
                     HashMap,Arrays,List,Map,Set,Exception,RuntimeException},
    keywordstyle=[2]\color{vscFunction},
    sensitive=true
}

\lstnewenvironment{java}[1][]{\lstset{style=javaStyle, #1}}{}
\newcommand{\javainline}[1]{\lstinline[style=javaStyle]!#1!}
\newcommand{\inputjava}[2][]{\lstinputlisting[style=javaStyle,#1]{#2}}       % Importa configuraciones para resaltar código

%---------------------------------------------------------------------------------------
% CONFIGURACIÓN DEL TEMA
%----------------------------------------------------------------------------------------
% Tema Base
\usetheme{Boadilla}                          % Define el tema principal
\useinnertheme{circles}                      % Tema interno con círculos
\useoutertheme{miniframes}                   % Tema externo con miniframes
\setbeamertemplate{navigation symbols}{}     % Elimina símbolos de navegación

% Colores Personalizados
\definecolor{primaryColor}{RGB}{11, 139, 171} % Color primario - cian Famaf oficial
\definecolor{secondaryColor}{RGB}{0, 104, 128}    % Color secundario - azul verdoso

% Configuraciones de Colores
\setbeamercolor{structure}{fg=primaryColor}
\setbeamercolor{palette primary}{bg=primaryColor, fg=white}
\setbeamercolor{palette secondary}{bg=secondaryColor, fg=white}
\setbeamercolor{title}{bg=primaryColor, fg=white}

% Colores del Encabezado y Pie de Página
\setbeamercolor{headline}{bg=secondaryColor, fg=white}
\setbeamercolor{section in head/foot}{bg=primaryColor, fg=white}
\setbeamercolor{subsection in head/foot}{bg=secondaryColor, fg=white}
\setbeamercolor{author in head/foot}{bg=primaryColor, fg=white}
\setbeamercolor{title in head/foot}{bg=secondaryColor, fg=white}
\setbeamercolor{date in head/foot}{bg=primaryColor, fg=white}
\setbeamercolor{page number in head/foot}{bg=primaryColor, fg=white}

%----------------------------------------------------------------------------------------
% BIBLIOGRAFÍA
%----------------------------------------------------------------------------------------
\bibliographystyle{plain}
%----------------------------------------------------------------------------------------
% INFORMACIÓN DE LA PRESENTACIÓN
%----------------------------------------------------------------------------------------
\title[Título]{Título completo}          % [Versión corta]{Versión completa}
\author[Nombre abreviado]{Nombre}        % [Versión corta]{Nombre completo}
\institute[FAMAF]{Facultad de Matemática, Astronomía, Física y Computación \\ (FAMAF - UNC) }
\date[Año]{MES / AÑO}

%----------------------------------------------------------------------------------------
% INICIO DEL DOCUMENTO
%----------------------------------------------------------------------------------------
\begin{document}

% Diapositiva de título con logo
\begin{frame}
    \begin{figure}
        \includegraphics[width=0.70\linewidth]{img/famaf.png}
    \end{figure}
    \titlepage
\end{frame}

% Sumario
\begin{frame}
    \frametitle{Estructura de la presentación}
    \tableofcontents
\end{frame}

% Inclusión de las secciones
\section{Ejemplos con texto}

%----------------------------------------------------------------------------------------

\begin{frame}
	\frametitle{Texto corrido}
    Lorem ipsum dolor sit amet, consectetur adipiscing elit. Nullam ipsum velit, cursus quis ligula eu, malesuada aliquet massa. Quisque non convallis felis, a auctor eros. Etiam sit amet turpis a sapien pulvinar malesuada quis quis nisi. Quisque scelerisque volutpat ligula vel mollis. Nam sit amet tristique erat, sit amet cursus mi. 
\end{frame}

%----------------------------------------------------------------------------------------

\begin{frame}
	\frametitle{Texto en tópicos numerados}
     Lorem ipsum dolor sit amet, consectetur adipiscing elit:
    \begin{enumerate}
        \item Lorem ipsum dolor sit amet.
        \item Lorem ipsum dolor sit amet.
    \end{enumerate}
	
\end{frame}

%----------------------------------------------------------------------------------------

\begin{frame}
	\frametitle{Texto en tópicos}
     Lorem ipsum dolor sit amet, consectetur adipiscing elit:
    \begin{itemize}
        \item Lorem ipsum dolor sit amet.
        \item Lorem ipsum dolor sit amet.
    \end{itemize}
	
\end{frame}



\section{Ejemplos con ecuaciones e imágenes}

%----------------------------------------------------------------------------------------

\begin{frame}{Una imagen}
   \begin{figure}[h]
       \centering
       \includegraphics[width=0.7\textwidth]{img/famaf.png}
       \caption{Leyenda}
       \label{fig:label_da_imagem}
   \end{figure}
\end{frame}

%----------------------------------------------------------------------------------------

\begin{frame}{Dos imágenes}
   \begin{figure}
       \centering
       \begin{subfigure}[b]{0.45\textwidth}
           \centering
           \includegraphics[width=\textwidth]{img/famaf.jpg}
           \caption{Leyenda 1}
           \label{fig:img1}
       \end{subfigure}
       \hfill
       \begin{subfigure}[b]{0.45\textwidth}
           \centering
           \includegraphics[width=\textwidth]{img/famaf.jpg}
           \caption{Leyenda 2}
           \label{fig:img2}
       \end{subfigure}
   \end{figure}
\end{frame}

%----------------------------------------------------------------------------------------
\begin{frame}{Ecuaciones}
    Ecuaciones de Navier-Stokes Forma expandida (3D):
    \footnotesize
        \begin{align*}
            \rho\left(\frac{\partial u}{\partial t} + u\frac{\partial u}{\partial x} + v\frac{\partial u}{\partial y} + w\frac{\partial u}{\partial z}\right) &= -\frac{\partial p}{\partial x} + \mu\left(\frac{\partial^2 u}{\partial x^2} + \frac{\partial^2 u}{\partial y^2} + \frac{\partial^2 u}{\partial z^2}\right) + f_x \\[0.3cm]
            \rho\left(\frac{\partial v}{\partial t} + u\frac{\partial v}{\partial x} + v\frac{\partial v}{\partial y} + w\frac{\partial v}{\partial z}\right) &= -\frac{\partial p}{\partial y} + \mu\left(\frac{\partial^2 v}{\partial x^2} + \frac{\partial^2 v}{\partial y^2} + \frac{\partial^2 v}{\partial z^2}\right) + f_y \\[0.3cm]
            \rho\left(\frac{\partial w}{\partial t} + u\frac{\partial w}{\partial x} + v\frac{\partial w}{\partial y} + w\frac{\partial w}{\partial z}\right) &= -\frac{\partial p}{\partial z} + \mu\left(\frac{\partial^2 w}{\partial x^2} + \frac{\partial^2 w}{\partial y^2} + \frac{\partial^2 w}{\partial z^2}\right) + f_z
        \end{align*}
        
    donde $\mathbf{v} = (u,v,w)$ es el campo de velocidad, $p$ es la presión, $\rho$ es la densidad, $\mu$ es la viscosidad dinámica y $\mathbf{f}$ representa fuerzas externas.
\end{frame}
\section{Ejemplos con código} 

%----------------------------------------------------------------------------------------
\begin{frame}[fragile]
    \frametitle{Python}
    
    \begin{python}
def calcular_doble(x):
    """Retorna el doble del número"""
    return 2 * x

# test de función
numero = 5
resultado = calcular_doble(numero)
print(f"El doble de {numero} es {resultado}")
    \end{python}
\end{frame}

%----------------------------------------------------------------------------------------
\begin{frame}[fragile]
    \frametitle{C}
    
    \begin{clang}
#include <stdio.h>

int main() {
    int numero = 5;
    int doble = 2 * numero;
    
    printf("el doble de %d es %d\n", numero, doble);
    return 0;
}
    \end{clang}
\end{frame}

%----------------------------------------------------------------------------------------
\begin{frame}[fragile]
    \frametitle{C++}
    
    \begin{cpp}
#include <iostream>
using namespace std;

int main() {
    int numero = 5;
    int doble = 2 * numero;
    
    cout << "el doble de " << numero;
    cout << " es " << doble << endl;
    return 0;
}
    \end{cpp}
\end{frame}

%----------------------------------------------------------------------------------------
\begin{frame}[fragile]
    \frametitle{R}
    
    \begin{rlang}
# Funcion para calcular el doble
calcular_doble <- function(x) {
  return(2 * x)
}

# Testeando la funcion
numero <- 5
resultado <- calcular_doble(numero)
print(paste("el doble de", numero, "es", resultado))
    \end{rlang}
\end{frame}

%----------------------------------------------------------------------------------------

\begin{frame}[fragile]
    \frametitle{Java}
    
    \begin{java}
public class Ejemplo {
    public static void main(String[] args) {
        int numero = 5;
        int doble = 2 * numero;
        
        System.out.println("el doble de " + numero +
                         " es " + doble);
    }
}
    \end{java}
\end{frame}
\section{Conclusión} 

\begin{frame}[allowframebreaks]{Referencias}
\nocite{*}
    \bibliography{bibliografia} 
\end{frame}

% Diapositiva final
\begin{frame}
    \begin{center}
        {\Huge ¡Fin de la presentación!}
    \end{center}
\end{frame}

\end{document}

